\scanpage{7}
- \mspage{4} -

utilise (1.1), en regardant les deux membres de (1.3) comme des sections de $p_*\mathcal O(6e)$ : en effet, $y^2-x^3$ est dans $p_*\mathcal O(5e)$, $xy$ dans $p_*\mathcal O(5e)$, $x^2$ dans $p_*\mathcal O(4e)$, $y$ dans $p_*\mathcal O(3e)$, $x$ dans $p_*\mathcal O(2e)$ et $1$ dans $p_*\mathcal O$.

Si on pose
\[
F(x,y)=y^2+a_1xy+a_3y-x^3-a_2x^2-a_4x-a_6
\]
la forme différentielle invariante définie par $\pi$ s'écrit
\[
\pi=dx/F_y=-dy/F_x\qquad\text{soit}
\]
\[
\pi=\frac{dx}{2y+a_1x+a_3}=\frac{dy}{3x^2+2a_2x+a_4-a_1y}.
\]
Réciproquement toute équation (1.3) est l'équation non homogène d'une courbe de genre 1 plongée comme une cubique dans $P_2$ et présentant un point d'inflexion en $(0,1,0)$ avec la droite de l'infini pour tangente d'inflexion.

On pose alors
\[
(1.4)\qquad b_2=a_1^2+4a_2\qquad b_4=a_1a_3+2a_4\qquad b_6=a_3^2+4a_6
\]
\[
-b_8=a_1a_3a_4+a_4^2-a_1^2a_6-a_2a_3^2-4a_2a_6
\]
\[
d_6=b_2b_4-18b_6\qquad d_8=b_2^2b_4-2^4\,3\,b_4^2+2\,3^2\,b_2b_6
\]
On a
\[
-4b_8=b_4^2-b_2b_6
\]
\[
12d_6=c_6+b_2c_4\qquad -12d_8=c_4^2+b_2c_6
\]
\margin{6}
